% Speaking script for presentation_pf_sssa_ts_gfm_en_v11.tex (16 slides).
%
% 15 pages: page 1 covers the title and contents slides together, then one page
% per remaining deck slide.  THREE OR FOUR LINES PER PAGE, at \large, so a page
% is taken in at a glance rather than read.  Budgets total 8:35, which is the
% point: the slot is fifteen minutes, so a slow delivery still fits and there is
% room for questions.
%
% RULE FOLLOWED THROUGHOUT: every fact spoken here is VISIBLE on the slide it
% belongs to.  A number readable only off a figure axis is spoken as "about".
\documentclass{beamer}
\usepackage{fontspec}
\setsansfont[Path=C:/Users/User/AppData/Local/Microsoft/Windows/Fonts/,UprightFont=Helvetica.ttf,BoldFont=Helvetica-Bold.ttf,ItalicFont=Helvetica-Oblique.ttf,BoldItalicFont=Helvetica-BoldOblique.ttf]{Helvetica.ttf}
\usetheme{Berlin}
\usefonttheme[onlymath]{serif}
\usepackage{xcolor}
\usepackage{amsmath}
\usepackage{bm}
\definecolor{myblue}{rgb}{0.0,0.0,0.7}
\definecolor{pnavy}{HTML}{0F2A4A}
\definecolor{pblue}{HTML}{1D5288}
\setbeamercolor{frametitle}{bg=myblue!82!black,fg=white}
\setbeamertemplate{navigation symbols}{}
\setbeamertemplate{background}{}
\setbeamertemplate{headline}{}
\setbeamertemplate{footline}{%
  \leavevmode%
  \hbox{%
  \begin{beamercolorbox}[wd=.86\paperwidth,ht=2.6ex,dp=1.1ex,leftskip=1.2em]{author in head/foot}%
    \scriptsize Speaking script --- An Automatic Following/Forming Framework for Inverter-Based Resources%
  \end{beamercolorbox}%
  \begin{beamercolorbox}[wd=.14\paperwidth,ht=2.6ex,dp=1.1ex,center]{date in head/foot}%
    \scriptsize \insertframenumber/\inserttotalframenumber
  \end{beamercolorbox}}%
}
\newcommand{\budget}[2]{%
  \vspace{-0.30cm}%
  \begin{flushright}\tiny\color{pnavy}%
  \textbf{#1}\quad$\vert$\quad total \textbf{#2}%
  \end{flushright}%
  \vspace{0.55cm}}
% \par first: without it the \vspace lands in horizontal mode and the cue runs
% on from the last sentence instead of standing on its own line.
\newcommand{\cue}[1]{%
  \par\vspace{0.55cm}{\scriptsize\color{pblue}\textit{Point at:} #1\par}}
% Beamer zeroes \parskip, which would run these short paragraphs together.
\setlength{\parskip}{0.55em}

\begin{document}
\begin{frame}{Script 1 --- Slides 1--2: Title and Contents}
\budget{0:25}{0:25}
\large
Good morning. I am Phichitchai Jueajan, with Pakkapol Phdungdan.

Our project is an automatic following and forming framework for inverter-based
resources.

In fifteen minutes: the problem, three methods, the switching law, four results.
\cue{the section list.}
\end{frame}

\begin{frame}{Script 2 --- Slide 3: Introduction}
\budget{0:30}{0:55}
\large
Left: grid-following needs an external voltage-angle reference.

Right: grid-forming makes its own angle, and can replace the reference lost when
the generator disconnects.

So neither mode is right all the time.
\cue{the last bullet on the right.}
\end{frame}

\begin{frame}{Script 3 --- Slide 4: Project Objectives}
\budget{0:30}{1:25}
\large
One in-house chain: power flow, small-signal stability, simulation. Four
converters with selectable branches.

Third objective is the heart of it: decide the allowed grid-forming sets from the
network condition, not by fixing the number in advance.
\cue{objective three.}
\end{frame}

\begin{frame}{Script 4 --- Slide 5: Overall Workflow}
\budget{0:35}{2:00}
\large
Top row: power flow, then equilibrium, then screening, then integration --- all
on one operating point.

Bottom row back: check severity, then make the transaction.

The rule: commit only after current continuity and network convergence.
\cue{the arrow back to the transaction box.}
\end{frame}

\begin{frame}{Script 5 --- Slide 6: Case Study}
\budget{0:25}{2:25}
\large
IEEE 14-bus: one generator, four dual-mode converters.

$\mathrm{SG}_1$ at bus one; $\mathrm{IBR}_1$ to $\mathrm{IBR}_4$ at buses two,
three, six and eight.

Every later figure is labelled by resource, not by bus.
\cue{the resource table.}
\end{frame}

\begin{frame}{Script 6 --- Slide 7: Power Flow Formulation}
\budget{0:35}{3:00}
\large
Net complex power at a bus: the voltage times the conjugate of the network
current. In polar form, the two real equations for $P$ and $Q$.

The table gives what each bus type fixes.

Below: Newton--Raphson --- mismatches, solve the boxed system, update, repeat.
\cue{the boxed Jacobian system.}
\end{frame}

\begin{frame}{Script 7 --- Slide 8: Small-Signal Stability}
\budget{0:35}{3:35}
\large
We linearise the same equations the simulator uses.

Since $\bm g_y$ is invertible, a Schur complement gives the boxed reduced matrix
$\bm A_{red}$.

Its eigenvalues give frequency and damping. Every physical root must have a
negative real part.
\cue{the two boxed equations.}
\end{frame}

\begin{frame}{Script 8 --- Slide 9: Time-Domain Simulation}
\budget{0:30}{4:05}
\large
Implicit trapezoidal for the states, with the network enforced at the new time
point --- solved together.

The list on the right is what makes an event trustworthy: land exactly on it,
apply the change atomically, require convergence after, roll back if it fails.
\cue{the event-handling block.}
\end{frame}

\begin{frame}{Script 9 --- Slide 10: Dual-Mode IBR Model}
\budget{0:35}{4:40}
\large
Each converter carries a fixed seventeen coordinates: shared plant, a
grid-following branch, a grid-forming branch.

The inactive branch is held, not reset --- that is why the mode can change
mid-run.

Right: the PLL, and the virtual synchronous machine.
\cue{the two active-set expressions.}
\end{frame}

\begin{frame}{Script 10 --- Slide 11: Switching Law}
\budget{0:40}{5:20}
\large
Severity $S$ is the average of a voltage index and a frequency index, saturated
to zero-to-one.

The command is deliberately asymmetric: into forming at $0.65$ held $0.10$
seconds --- fast; back to following below $0.35$ held a full second --- slow.

And a gate: the current must match to $10^{-10}$ per unit.
\cue{the two dwell arrows.}
\end{frame}

\begin{frame}{Script 11 --- Slide 12: Result 1, Candidate Screening}
\budget{0:40}{6:00}
\large
Any single converter passes. Two pairs pass, and a pair gives the best margin.
Every triple is rejected. All four together passes again.

Three fails and four passes --- admissibility is non-monotonic in the count.

So every candidate must be computed, not inferred.
\cue{the three-converter row, then the four-converter row.}
\end{frame}

\begin{frame}{Script 12 --- Slide 13: Result 2, Disturbance Scenario}
\budget{0:35}{6:35}
\large
Six periods over $250$ seconds: the generator trips at $20$, load steps at $50$,
fault at $85$, line trip at $110$, reclose commanded at $145$.

The two times under period six are results, read from the event log --- not
settings.
\cue{the two result times under period six.}
\end{frame}

\begin{frame}{Script 13 --- Slide 14: Result 3, Electrical Response}
\budget{0:40}{7:15}
\large
Panels (a) and (b): the converters pick up the island after $20$ seconds, then
hand it back.

Panel (c): each disturbance drives severity above $\Gamma_{on}$, and between
events it falls below $\Gamma_{off}$ --- so this is measurement, not a timetable.

Panel (d): the reference moves to $\mathrm{IBR}_1$, then back.
\cue{a peak in (c), then the matching step in (d).}
\end{frame}

\begin{frame}{Script 14 --- Slide 15: Result 4, Bus 9 Comparison}
\budget{0:40}{7:55}
\large
Blue is adaptive, grey is fixed following. Bus nine carries load but no device,
so it shows what the island delivers.

In (a) to (c) blue holds near one per unit; grey sags to about two-thirds and
delivers less.

In (d), the honest trade-off: grey holds a narrower frequency band.
\cue{the gap in (a), then the band in (d).}
\end{frame}

\begin{frame}{Script 15 --- Slide 16: Conclusion}
\budget{0:40}{8:35}
\large
Three findings: after the generator disconnects at least one forming converter is
required; admissibility is non-monotonic; staged switching reaches $250$ seconds
and hands everything back.

The box compares the arms --- locked following stops at $20$ seconds, four pinned
units fail at $25.485$.

Thank you --- we welcome questions.
\cue{the comparative-outcome box.}
\end{frame}


\end{document}
